\documentclass[10pt,a4paper]{article}

\usepackage[margin=0.85in]{geometry}
\usepackage[utf8]{inputenc}
\usepackage[T1]{fontenc}
\usepackage{lmodern}
\usepackage{microtype}
\usepackage{amsmath}
\usepackage{amssymb}
\usepackage{booktabs}
\usepackage{tabularx}
\usepackage{array}
\usepackage{enumitem}
\usepackage{xcolor}
\usepackage{hyperref}
\usepackage{textcomp}

\hypersetup{colorlinks=true, linkcolor=black, urlcolor=blue}
\setlist[itemize]{itemsep=1pt, topsep=2pt, parsep=0pt}
\setlist[enumerate]{itemsep=1pt, topsep=2pt, parsep=0pt}

\newcommand{\code}[1]{\texttt{\small #1}}

\title{\textbf{Automatic Creation of Indices in PostgreSQL\\\large Project Submission Report}}
\author{}
\date{}

\begin{document}

\maketitle
\vspace{-3.5em}

\begin{center}
\begin{tabular}{ll}
\toprule
\textbf{Name} & \textbf{Roll number} \\
\midrule
Siddhant Mulkikar         & 23B0956 \\
Yuvraj Gupta              & 23B0999 \\
Sandeep Reddy Nallamilli  & 23B1006 \\
Yash Sabale               & 23B1043 \\
\bottomrule
\end{tabular}
\end{center}

\vspace{0.5em}
\noindent\textbf{Public git repository (forked from \code{postgres/postgres}):}\\
\url{https://github.com/HypedUpbOii/postgres}

\vspace{0.5em}
\noindent The repository is a fork of the upstream PostgreSQL source tree.
The initial pull from upstream is preserved as a single ancestor commit;
all our work is on top of it under \code{contrib/auto\_index/},
\code{scripts/}, and \code{report/}. Per-file diffs against upstream are
included in the submission zip under \code{diffs/}.

\section{Motivation}

Indexes are the single biggest lever a relational database has on read
latency --- turning an \(O(N)\) filter into an \(O(\log N)\) lookup ---
but they are notoriously hard to manage manually. Workloads change over
time, indexes have ongoing write-amplification and storage cost,
composite-key ordering is non-obvious (\code{(a,\,b)} or
\code{(b,\,a)}?), and most teams add indexes reactively and never
remove them. Oracle 19c ships an ``Auto Indexing'' feature; SQL Server
has Auto Tuning. PostgreSQL has \code{hypopg} (hypothetical indexes)
and \code{pg\_qualstats} (predicate counters), but \emph{no autonomous
create/drop loop} --- a DBA still has to interpret recommendations and
act. Our project closes that gap: a PostgreSQL extension,
\code{auto\_index}, that observes the live workload and creates and
drops indexes (single-column or composite, up to three columns)
entirely on its own.

\section{Functionality implemented}

The extension hooks into PostgreSQL in three places. An
\code{ExecutorEnd} hook runs after every query, walks the plan tree,
finds every \code{SeqScanState}, and decodes its filter quals into
\code{(attno, op)} pairs. The decoder handles \code{Var op Const},
\code{IN (\dots)}, \code{BETWEEN}, AND/OR via \code{BoolExpr} (so even
TPC-H Q19's deeply nested OR-of-AND extracts correctly), and implicit
casts via \code{RelabelType}. A fixed-size shared-memory area
(\code{AutoIndexSharedState}) holds per-table predicate counters: per
column ``singleton'' counts and a separate ``colset'' bucket that
records when 2--3 columns co-occur in the same query. A background
worker wakes every \code{check\_interval} seconds, folds interval
counters into cumulative state, runs the active create-strategy and
the drop logic, and emits \code{CREATE INDEX} / \code{DROP INDEX}
directly via SPI.

\paragraph{Composite indexes (up to 3 columns).} When the same group
of 2--3 columns repeatedly co-occurs, the worker builds a composite.
Column order at create time puts equality columns first, range columns
last (B-tree cannot use later columns past a range), with attno as the
tiebreak. To avoid redundant work, the singleton index on the
composite's \emph{leading} column is suppressed (subsumed by the
composite's leftmost prefix), but singletons on non-leading columns
are kept because B-tree cannot use them on their own.

\paragraph{Pluggable create-strategies.} Six strategies are selectable
through the \code{auto\_index.create\_strategy} GUC:

\begin{center}\small
\begin{tabular}{ll}
\toprule
\textbf{Strategy} & \textbf{Rule} \\
\midrule
\code{ski\_rental} \emph{(default)} & \code{cumulative\_benefit} $\geq \lceil$\code{build\_cost}/\code{seq\_scan\_cost}$\rceil$, AND read/write ratio above threshold \\
\code{simple\_threshold}            & \code{cumulative\_benefit} $\geq N$ \\
\code{ratio\_only}                  & read/write ratio above threshold \\
\code{cost\_gain}                   & projected savings minus maintenance must pay off \code{build\_cost} \\
\code{size\_gated}                  & \code{simple\_threshold} + minimum table size \\
\code{always}                       & fire on any positive benefit \\
\bottomrule
\end{tabular}
\end{center}

Each strategy is a single \code{decide\_create\_*(ctx) $\to$ bool}
function; the dispatcher is one switch.

\paragraph{Drop logic.} Drop is ski-rental in reverse. For every
auto-created index (we keep our own catalog table that survives
restart) the worker reads \code{pg\_stat\_user\_indexes.idx\_scan}
each interval. If \code{idx\_scan} increased, the idle accumulator is
reset; otherwise \code{(writes $\times$ maint\_per\_write +
seq\_scan\_cost) $\times$ 1000} is added. Once the accumulator passes
\code{build\_cost $\times$ 1000} the index is dropped. The
\code{seq\_scan\_cost} floor ensures even zero-write tables eventually
age out an index that is never read.

\paragraph{SQL surface.} \code{auto\_index\_reset()},
\code{auto\_index\_stats()}, and the catalog view
\code{auto\_index\_catalog} (with the helper
\code{auto\_index\_attnames(relid, attnos)} for human-readable column
names).

\paragraph{Scope limitations} (documented in code, not
security-critical): predicate extraction only fires on
\code{SeqScanState} --- once any other index exists on a table, the
planner switches to \code{Bitmap Heap Scan} and our hook can no longer
see filter predicates (the dominant case in OLTP, capping TPC-C gains
around 1.14$\times$). Aggregate-bound queries (TPC-H Q01 doing
\code{SUM} over the whole table) cannot be helped by any index in
principle.

\section{Code created}

All new code lives under \code{contrib/auto\_index/} (the extension)
and \code{scripts/} (build and benchmark drivers). \emph{No upstream
PostgreSQL source files were modified} --- the extension lives
entirely in \code{contrib/}, hooks into existing extension points, and
links cleanly against an unmodified server.

\begin{center}\small
\begin{tabular}{lll}
\toprule
\textbf{Path} & \textbf{Purpose} & \textbf{Lines} \\
\midrule
\code{contrib/auto\_index/auto\_index.c}        & core: hooks, shmem, bgworker, parsing, strategies, DDL emission & 1\,802 \\
\code{contrib/auto\_index/auto\_index.h}        & shared structs and GUCs declaration & 68 \\
\code{contrib/auto\_index/auto\_index--1.0.sql} & SQL functions and catalog table & 37 \\
\code{contrib/auto\_index/auto\_index.control}, \code{Makefile} & extension manifest, PGXS build & --- \\
\code{scripts/\{build,run,benchmark,demo\}.sh}  & top-level driver scripts & $\sim$250 \\
\code{scripts/lib/*.sh}                         & per-benchmark helpers (pgbench, TPC-C, TPC-H, strategies, setup) & $\sim$700 \\
\code{scripts/tpch\_queries/q*.sql}             & 8 representative TPC-H queries & --- \\
\code{scripts/tpcc\_xn/*.sql}, schema/load      & TPC-C-lite schema and the four transactions & --- \\
\code{Dockerfile}, \code{docker-compose.yaml}   & reproducible build container & --- \\
\bottomrule
\end{tabular}
\end{center}

The C source is split into six clearly-labelled sections inside
\code{auto\_index.c}: \code{1.\,MODULE INIT \& GUCs},
\code{2.\,SHARED MEMORY}, \code{3.\,PARSING}, \code{4.\,STRATEGIES},
\code{5.\,INDEX MANAGEMENT}, and \code{6.\,SQL-CALLABLE FUNCTIONS}.

\section{How we tested}

Four benchmarks, scripted under \code{scripts/benchmark.sh}. All
numbers come from a Docker container (Ubuntu 22.04, 6 cores, 8\,GB
RAM, NVMe SSD), PostgreSQL 19devel built with \code{-O2
--enable-debug --enable-cassert}. Each scenario was run three times.

\paragraph{pgbench live demo.} 300\,K-row table, 4 clients $\times$
2 worker threads, single equality predicate; TPS reported every
2\,s.

\begin{table}[h]
\centering\small
\begin{tabular}{lrrrl}
\toprule
\textbf{Window} & \textbf{Run 1} & \textbf{Run 2} & \textbf{Run 3} & \textbf{Phase} \\
\midrule
$t=2$~s        & 380     & 342     & 316     & seq-scan baseline \\
$t=4$~s        & 340     & 337     & 330     & seq-scan baseline \\
$t=6$~s        & \textbf{14\,696} & \textbf{15\,116} & \textbf{15\,499} & \textbf{auto\_index fires} \\
$t=8$~s        & 17\,571 & 18\,080 & 17\,859 & indexed (steady) \\
$t=10$--18~s   & 17\,200--18\,400 & 17\,650--18\,460 & 17\,500--18\,400 & indexed (steady) \\
\bottomrule
\end{tabular}
\end{table}

The step happens at $t=6$~s in every run --- the bgworker's second
snapshot. Steady-state TPS is roughly 17\,800 against a baseline of
$\sim$340: a factor of $\sim$52 in throughput, with no human asking
for an index.

\paragraph{TPC-H, all six strategies.} Same TPC-H workload (8
representative queries: Q1, Q3, Q5, Q6, Q10, Q12, Q14, Q19) under
each of the six strategies. The same baseline phase was used as a
reference for all strategies; for each strategy we restarted the
server, ran 5 training passes, waited 20\,s, and measured.

\begin{table}[h]
\centering\small
\begin{tabular}{lrrrrrrr}
\toprule
\textbf{Query} & \textbf{Baseline} & \multicolumn{6}{c}{\textbf{Speedup vs.\ baseline (higher is better)}} \\
\cmidrule(lr){3-8}
                &  \textbf{(ms)} & \textbf{ski\_rental} & \textbf{simple\_thr} & \textbf{ratio\_only} & \textbf{cost\_gain} & \textbf{size\_gated} & \textbf{always} \\
\midrule
q01   & 4\,916  & 1.05$\times$ & 1.04$\times$ & 1.04$\times$ & 1.02$\times$ & 1.02$\times$ & 0.99$\times$ \\
q03   &   978   & 1.07$\times$ & 0.99$\times$ & 1.02$\times$ & 1.10$\times$ & 0.98$\times$ & 0.96$\times$ \\
q05   &   650   & 1.48$\times$ & 0.94$\times$ & 1.60$\times$ & 1.64$\times$ & 0.55$\times$ & 1.05$\times$ \\
q06   &   601   & 1.57$\times$ & 2.61$\times$ & \textbf{3.20$\times$} & 1.50$\times$ & 1.65$\times$ & 2.62$\times$ \\
q10   & 1\,016  & 1.28$\times$ & 1.22$\times$ & 1.24$\times$ & 1.30$\times$ & 1.12$\times$ & 1.24$\times$ \\
q12   & 1\,011  & 1.47$\times$ & 1.96$\times$ & \textbf{2.06$\times$} & 1.48$\times$ & 1.75$\times$ & 1.94$\times$ \\
q14   &   551   & 2.80$\times$ & 2.52$\times$ & 2.53$\times$ & \textbf{2.91$\times$} & 2.51$\times$ & 2.49$\times$ \\
q19   &   855   & 2.13$\times$ & 2.82$\times$ & \textbf{3.07$\times$} & 2.18$\times$ & 2.78$\times$ & 2.93$\times$ \\
\midrule
\textbf{geomean}       & --- & 1.53$\times$ & 1.60$\times$ & \textbf{1.79$\times$} & 1.55$\times$ & 1.32$\times$ & 1.61$\times$ \\
\textbf{indexes built} & --- & 6 & 11 & 10 & 6 & 10 & 13 \\
\bottomrule
\end{tabular}
\end{table}

\noindent\code{ratio\_only} comes out on top with a geomean of
$\sim$1.79$\times$, achieving 3$\times$+ speedups on Q06, Q12, and
Q19 by building the right composite indexes (notably the 3-column
\code{(l\_quantity, l\_discount, l\_shipdate)} that exactly matches
Q06's filter). \code{ski\_rental}, the conservative default, still
gets a respectable 1.53$\times$ with only 6 indexes.

\paragraph{TPC-C-lite (OLTP).} 2 warehouses, pgbench-driven mix of
NewOrder / Payment / OrderStatus / StockLevel. The single
auto-created index was on \code{orders.o\_c\_id}; \code{auto\_index}
correctly \emph{avoided} indexing customer columns despite heavy
reads, because Payment also writes to those rows. Mean speedup
1.14$\times$ $\pm$ 0.02$\times$ across runs --- a small but correct
OLTP result, capped by the bitmap-heap-scan limitation noted above.

\paragraph{Lifecycle demo.} 50\,K-row table; phase A is 200 reads on
a single column then 200 reads on a 3-column predicate; phase B is
300 INSERT/UPDATE/DELETE rounds with \emph{no} reads on the indexed
columns.

\begin{table}[h]
\centering\small
\begin{tabular}{crlc}
\toprule
\textbf{Run} & \textbf{Phase A: indexes built} & \textbf{Phase B: time to drop all} & \textbf{Result} \\
\midrule
1 & 4 (singletons + composite) & $\sim$20~s & \checkmark \\
2 & 4 (singletons + composite) & $\sim$20~s & \checkmark \\
3 & 4 (singletons + composite) & $\sim$20~s & \checkmark \\
\bottomrule
\end{tabular}
\end{table}

Every run produced both a singleton and a composite during phase A,
and then dropped all of them roughly four bgworker cycles into
phase B. The drop accumulator is deterministic enough that the
$\sim$20\,s figure is practically a constant.

A late-stage docker-compose fix: PostgreSQL parallel hash joins back
their shared hash tables with POSIX shm under \code{/dev/shm}, and
Docker's 64\,MB default caused TPC-H Q10/Q12/Q14 to fail with
``could not resize shared memory segment''. We set
\code{shm\_size:\,1g} in \code{docker-compose.yaml} and added a
named volume for \code{/usr/local/pgsql} so future compose-file
edits don't nuke the install.

\section{How we used AI}

AI was \emph{not} used to design the high-level approach (the
ski-rental analogy, the strategy-pluggable architecture, the choice
of hook points); those came from team discussion. The line-by-line
C, the qual-extraction code, the benchmark plumbing, and most of
the internals research were AI-assisted.

\textbf{We have read and verified every line of code that AI
produced.} Each chunk of AI-generated C went through compile + run +
regression of an earlier benchmark before being committed; several
AI-suggested approaches were rejected for being non-idiomatic
(e.g.\ an early draft used \code{palloc} in shared memory; we
rewrote with \code{ShmemAlloc}).

The kinds of prompts we used:

\begin{itemize}
\item Locating where filter quals live in the executor at runtime.
\item The standard chain-and-restore idiom for an \code{ExecutorEnd\_hook}.
\item \code{shmem\_request\_hook} + \code{shmem\_startup\_hook} pattern.
\item A recursive walker over \code{OpExpr} / \code{BoolExpr} /
      \code{ScalarArrayOpExpr} that peels \code{RelabelType}.
\item Emitting \code{CREATE INDEX} from a bgworker safely (SPI).
\item Which prefixes a multi-column B-tree can satisfy.
\item An extension-owned catalog table that survives restart.
\item Bash plumbing for the per-strategy comparison table.
\item Debugging the Docker \code{/dev/shm} parallel-hash failure.
\item Drafting prose from team-supplied bullet points and benchmark CSVs.
\end{itemize}

We have checked every line AI has written for correctness.

\section{Submission zip layout}

\begin{verbatim}
auto_index_submission/
|-- README.md                       # build & run instructions
|-- git_url.txt                     # link to public fork
|-- diffs/                          # per-file git diffs vs. upstream PostgreSQL
|   |-- .gitignore.diff             # the only modified upstream file
|-- new_files/                      # all new files we authored, mirrored at original paths
|   |-- contrib/auto_index/...
|   |-- scripts/...
|   |-- Dockerfile, docker-compose.yaml
|-- test_data/                      # test artefacts
|   |-- tpch_queries/q*.sql, tpch_schema.sql
|   |-- tpcc_schema.sql, tpcc_load.sql, tpcc_xn/*.sql
|-- report/
|   |-- submission_report.pdf       # this document
|   |-- report.pdf                  # full technical report (deeper dive)
\end{verbatim}

The TPC-H raw data files (\code{*.tbl}, $\sim$70\,MB at SF=0.1) are
not shipped --- \code{scripts/lib/tpch\_setup.sh} regenerates them
with \code{tpch-dbgen} on first run.

\paragraph{Reproducing the headline results from a clean checkout.}

\begin{verbatim}
git clone https://github.com/HypedUpbOii/postgres.git
cd postgres
scripts/build.sh all                # docker + pg + init + extension
scripts/benchmark.sh pgbench        # ~52x live step
scripts/benchmark.sh tpch           # 8-query 3-phase TPC-H
scripts/benchmark.sh tpcc           # OLTP
scripts/demo.sh                     # full create-then-drop lifecycle, ~90 s
\end{verbatim}

\end{document}
